\chapter{Analyse de l'Existant et Spécification des Besoins}

\section{Analyse Critique des Systèmes de Supervision Conventionnels}

Dans la plupart des environnements industriels lourds, la supervision repose historiquement sur des systèmes de contrôle et d'acquisition de données (\textbf{SCADA --- Supervisory Control and Data Acquisition}). Si ces systèmes s'avèrent fiables pour le pilotage d'automatismes locaux (automates programmables industriels --- API/PLC), ils révèlent plusieurs limitations majeures face aux exigences modernes de maintenance préventive et d'aide à la décision :

\begin{itemize}
    \item \textbf{Cloisonnement des mesures :} Chaque capteur est supervisé de manière unitaire. Le système déclenche une alarme lorsqu'une valeur franchit un seuil statique, sans mise en relation contextuelle avec l'état des machines situées en amont ou en aval dans la chaîne de production ;
    \item \textbf{Phénomène de « fatigue d'alarme » (\textit{Alarm Fatigue}) :} Lors d'une dérive de fonctionnement, un capteur oscillant autour d'une valeur limite génère des dizaines de notifications identiques par minute. Cette surcharge cognitive désensibilise l'opérateur et peut masquer des incidents critiques ;
    \item \textbf{Incapacité à projeter l'impact en cascade :} Les systèmes conventionnels n'intègrent pas la topologie du réseau de production. Lorsque le concasseur principal s'arrête brutalement, l'opérateur doit déduire mentalement les équipements avals qui vont subir un étranglement d'alimentation ;
    \item \textbf{Rigidité et dette technique des interfaces :} Les interfaces homme-machine (IHM) traditionnelles sont souvent statiques, conçues avec des technologies propriétaires peu évolutives et dépourvues d'interactivité topologique moderne.
\end{itemize}

\section{Étude Comparative : SCADA vs Jumeau Numérique}

Le Tableau~\ref{tab:comparaison_scada_dt} met en évidence les gains qualitatifs et fonctionnels apportés par le paradigme du Jumeau Numérique face aux solutions de supervision conventionnelles.

\begin{table}[htbp]
\centering
\small
\caption{Comparatif entre la supervision SCADA classique et l'approche Jumeau Numérique}
\label{tab:comparaison_scada_dt}
\begin{tabularx}{\textwidth}{p{3.5cm} Y Y}
\toprule
\textbf{Critère} & \textbf{Supervision SCADA Classique} & \textbf{Jumeau Numérique Développé} \\
\midrule
\textbf{Modélisation} & Table de registres d'automates & Entités orientées domaine métier (DDD) \\
\textbf{Gestion des alertes} & Seuils statiques, alertes répétitives & Évaluation multi-seuils avec cooldown anti-rebond \\
\textbf{Propagation d'impact} & Absente (déduction manuelle) & Algorithmique (Parcours BFS sur graphe topologique) \\
\textbf{Architecture logicielle} & Monolithique propriétaire & Clean Architecture modulaire, CQRS et MediatR \\
\textbf{Diffusion temps réel} & Rafraîchissement périodique (Polling) & Poussée instantanée WebSockets (SignalR) \\
\textbf{Expérience utilisateur} & Synoptiques 2D statiques & Tableau de bord réactif Angular 19 (vis-network, Chart.js) \\
\bottomrule
\end{tabularx}
\end{table}

\section{Spécification des Besoins Fonctionnels}

Les fonctionnalités du système ont été organisées autour de cinq axes majeurs :

\begin{enumerate}[label=\textbf{BF.\arabic* ---}, leftmargin=2.2cm]
    \item \textbf{Ingestion continue de télémétrie :} Réception des flux de mesures simulées (vibration, température, pression, courant, débit, niveau) horodatées à haute fréquence.
    \item \textbf{Évaluation dynamique des anomalies :} Analyse de chaque mesure entrante par rapport aux seuils configurés sur le capteur (\textit{Warning} et \textit{Critical}) pour classifier son niveau de conformité.
    \item \textbf{Régulation anti-rebond (\textit{Alert Cooldown}) :} Mécanisme de temporisation glissante empêchant l'émission répétée d'alertes identiques pour un même capteur durant un intervalle paramétrable (ex. 60 secondes).
    \item \textbf{Analyse d'impact en cascade :} En cas de franchissement d'un seuil critique sur un équipement source, exécution automatique d'un algorithme de parcours en largeur (BFS) sur le graphe de dépendance de l'usine pour identifier l'ensemble des machines avals menacées.
    \item \textbf{Gestion du cycle de vie des alertes :} Traçabilité complète des alertes avec statuts (\textit{Active}, \textit{Acknowledged}, \textit{Resolved}) et possibilité pour l'opérateur d'acquitter une alarme avec horodatage.
    \item \textbf{Cartographie topologique interactive :} Rendu graphique dynamique du réseau d'équipements sous forme de graphe interactif coloré selon l'état de santé opérationnel instantané des machines.
    \item \textbf{Visualisation temporelle de la télémétrie :} Graphiques dynamiques des séries chronologiques des capteurs pour visualiser les tendances de dégradation.
    \item \textbf{Moteur de simulation et rejeu d'incidents :} Générateur de flux synthétiques autonome capable d'exécuter des scénarios de pannes déterministes pour la validation du système.
\end{enumerate}

\section{Spécification des Besoins Non-Fonctionnels}

Le système doit satisfaire les exigences techniques et de qualité logicielle suivantes :

\begin{itemize}
    \item \textbf{Performance et Latence (BNF-1) :} Le traitement complet d'une mesure dans le pipeline backend (validation, analyse d'anomalie, propagation BFS et notification SignalR) doit être inférieur à $50$\,ms.
    \item \textbf{Intégrité et Invariants Métier (BNF-2) :} Les règles métier (ex. transitions d'état d'une machine, validité des seuils de capteurs, dépendances acycliques) doivent être strictement encapsulées et protégées contre les états invalides via le \textit{DomainGuard}.
    \item \textbf{Extensibilité et Découplage (BNF-3) :} L'architecture logicielle doit respecter l'inversion des dépendances pour permettre l'évolution future des connecteurs de données (ex. passage d'un simulateur à un broker MQTT/OPC-UA industriel) sans altérer le cœur de domaine.
    \item \textbf{Ergonomie et Réactivité de l'IHM (BNF-4) :} L'interface Web doit être réactive, fluide (60 FPS) et permettre une prise en main immédiate par les équipes d'exploitation grâce à un code couleur normalisé.
    \item \textbf{Couverture de Tests (BNF-5) :} Le cœur métier et les algorithmes critiques doivent faire l'objet d'une suite de tests unitaires automatisés exhaustive.
\end{itemize}
